\section{Exercise 7}

\begin{frame}{Information Gain Ratio and Classifiers}
\textbf{Description}

In this exercise we create a classifier, implemented as a decision tree,
using the principle of maximum information gain ratio. A decision tree classifier is an algorithm that classifies a vector of feaures using the following iterative steps for each node of the tree

\begin{itemize}
	
	\item If all the vector in the training set belong to the same class
	
	\item For each node, it analyzes a single feature and moves on the next node based on the result of the analysis. If the feature has limited number of possibile values, a branch is created for each value. If the feature has a numerical value, the value is compared with a threshold and a branch is created for each result of the comparison
	
	\item If all the vector in the training set 



\end{itemize}

analyzing a single feature in each node and moving on the next node based on the result of the analysis, until a leaf

This principle states that the features

\end{frame}

The principle of maximum 


\begin{frame}
	
\end{frame}
